% Part: first-order-logic
% Chapter: proof-systems
% Section: seqeunt-calculus

\documentclass[../../../include/open-logic-section]{subfiles}

\begin{document}

\iftag{FOL}
      {\olfileid{fol}{prf}{seq}}
      {\olfileid{pl}{prf}{seq}}

\olsection{সিকোয়েন্ট কলন}

অনেক !!{derivation} পদ্ধতি !!{sentence}s-এর বিন্যাস নিয়ে কাজ করলেও,
সিকোয়েন্ট কলন \emph{সিকোয়েন্ট} নিয়ে কাজ করে। সিকোয়েন্ট হলো নিচের
রূপের অভিব্যক্তি:
\[
!A_1, \dots, !A_m \Sequent !B_1, \dots, !B_n,
\]
অর্থাৎ সিকোয়েন্ট সংকেত~$\Sequent$ দিয়ে পৃথক করা !!{sentence}s-এর
দুটি অনুক্রমের একটি যুগল। অনুক্রমদুটির যেকোনোটিই শূন্য হতে পারে।
% BN-SRC-021 TARGET-CORRECTION-BEGIN
\par\smallskip\noindent\textnormal{\small\textbf{উৎস-সংশোধন (BN-SRC-021):} হিমায়িত ইংরেজি উৎসে সিকোয়েন্টের দুপাশের শেষ সূচকই~$m$, যদিও পরের বাক্যেই দুই অনুক্রম স্বাধীনভাবে শূন্য হতে পারে বলা হয়েছে। বাংলা পাঠে তাদের স্বাধীন দৈর্ঘ্য বোঝাতে ডানপাশের শেষ সূচক~$n$ করা হয়েছে।} % BN-SRC-021 TARGET-CORRECTION-END
সিকোয়েন্ট কলনে কোনো !!^a{derivation} হলো সিকোয়েন্টের একটি বৃক্ষ,
যেখানে সর্বোচ্চ সিকোয়েন্টগুলি বিশেষ রূপের (তাদের “প্রারম্ভিক সিকোয়েন্ট”
বা “স্বতঃসিদ্ধ” বলা হয়), আর অন্য প্রতিটি সিকোয়েন্ট তার ঠিক উপরের
সিকোয়েন্টগুলি থেকে কোনো একটি অনুমান-বিধি অনুসারে আসে। অনুমান-বিধিগুলি
হয় সিকোয়েন্টের !!{sentence}s-কে বদলায় (বাঁদিকে বা ডানদিকে যোগ করে,
বাদ দেয় কিংবা পুনর্বিন্যস্ত করে), নয়তো বিধির সিদ্ধান্তে একটি জটিল
!!{formula} প্রবর্তন করে। যেমন, $\LeftR{\land}$ বিধি $!A,
\Gamma \Sequent \Delta$ থেকে $!A \land !B, \Gamma \Sequent \Delta$
অনুমান করতে দেয়, আর $\RightR{\lif}$ বিধি $!A, \Gamma \Sequent
\Delta, !B$ থেকে $\Gamma \Sequent \Delta, !A \lif !B$ অনুমান করতে
দেয়—যেকোনো $\Gamma$, $\Delta$, $!A$ ও~$!B$-এর জন্য। (বিশেষত,
$\Gamma$ ও~$\Delta$ শূন্য হতে পারে।)

সিকোয়েন্ট কলনভিত্তিক $\Proves$ সম্পর্কটি নিচের মতো সংজ্ঞায়িত:
$\Gamma \Proves !A$ তখনই, যখন এমন কোনো অনুক্রম $\Gamma_0$ আছে যাতে
$\Gamma_0$-এর প্রতিটি $!A$,~$\Gamma$-র অন্তর্গত এবং এমন একটি
!!{derivation} আছে যার মূলে~$\Gamma_0 \Sequent !A$ সিকোয়েন্টটি রয়েছে।
সিকোয়েন্ট কলনে $!A$ একটি উপপাদ্য, যদি~$\Sequent !A$ সিকোয়েন্টটির
একটি !!a{derivation} থাকে। যেমন, নিচের !!a{derivation}-টি দেখায় যে
$\Proves (!A \land !B) \lif !A$:
\begin{prooftree}
  \Axiom$!A \fCenter !A$
  \RightLabel{\LeftR{\land}}
  \UnaryInf$!A \land !B \fCenter !A$
  \RightLabel{\RightR{\lif}}
  \UnaryInf$\fCenter (!A \land !B) \lif !A$
\end{prooftree}

সিকোয়েন্ট কলনে কোনো সেট $\Gamma$ অসঙ্গত, যদি $\Gamma_0 \Sequent$-এর
একটি !!a{derivation} থাকে (যেখানে প্রতিটি $!A \in \Gamma_0$,
~$\Gamma$-র অন্তর্গত এবং সিকোয়েন্টের ডানপাশ শূন্য)।
\RightR{\Weakening} বিধি ব্যবহার করে কোনো অসঙ্গত সেট থেকে যেকোনো
!!{sentence} !!{derive}d করা যায়।

গেরহার্ড গেন্‌ৎসেন ১৯৩০-এর দশকে সিকোয়েন্ট কলন উদ্ভাবন করেন। এর নিয়মবদ্ধ
ও প্রতিসম নকশার কারণে !!{derivation}s-এর তত্ত্ব গড়ে তুলতে এটি অত্যন্ত
উপযোগী একটি রীতিব্যবস্থা। সিকোয়েন্ট কলনে !!{derivation}s খুঁজে পাওয়া
তুলনামূলকভাবে সহজ, কিন্তু এই !!{derivation}s প্রায়ই পড়া কঠিন এবং
প্রমাণের সঙ্গে তাদের সম্পর্কও কখনো কখনো সহজে ধরা পড়ে না। তবু
!!{derivation} পদ্ধতির ক্ষেত্রে এটি অত্যন্ত সুন্দর একটি উপায় হিসেবে
প্রমাণিত হয়েছে, এবং বহু যুক্তিবিদ্যার জন্য সিকোয়েন্ট কলন পদ্ধতি আছে।

\end{document}
